\begin{center}
\begin{tikzpicture}[scale=1.2,>=latex]
% -------------------------------------- Defino estilo de flecha -----------------------------------
\tikzset{
    mid arrow/.style={
        decoration={
            markings,
            mark=at position 0.5 with {\arrow{>}} 
        },
        postaction={decorate} 
    },
    init arrow/.style={
        decoration={
            markings,
            mark=at position 0.2 with {\arrow{>}} 
        },
        postaction={decorate} 
    }
}
\begin{axis}[
    hide axis, clip=false,
    view={90}{40}
]
%------------------------------------- Gráfica Banda de Mobius -------------------------------------
\addplot3 [
    surf, shader=faceted interp,
    point meta=x,
    % color=orange,% Color sólido
    colormap={graywhite}{color=(gray!70!white) color=(white)}, % Gradiente suave
    samples=40,
    samples y=5,
    z buffer=sort,
    domain=0:360,
    y domain=-0.5:0.5
] (
    {(1+0.5*y*cos(x/2)))*cos(x)},
    {(1+0.5*y*cos(x/2)))*sin(x)},
    {0.5*y*sin(x/2)});

%-------------------------------------- Primer segmento frontera e_4 --------------------------------------------
\addplot3 [
    samples=50,
    samples y=0, % Evita saltos abruptos en y
    domain=180:387,       
    thick,     
    red,       
    init arrow             
] (
    {(1+0.25*cos(x/2)))*cos(x)},  % Punto en el borde (y = 0.5)
    {(1+0.25*cos(x/2)))*sin(x)},  
    {0.25*sin(x/2)}        
) node[pos=0.2, above] {$e_4$}; % <- El nombre se añade aquí
\addplot3 [
    samples=50,
    samples y=0, % Evita saltos abruptos en y
    domain=396:450,       
    thick,     
    red,                    
] (
    {(1+0.25*cos(x/2)))*cos(x)},  % Punto en el borde (y = 0.5)
    {(1+0.25*cos(x/2)))*sin(x)},  
    {0.25*sin(x/2)}        
);
\addplot3 [
    samples=50,
    samples y=0, % Evita saltos abruptos en y
    domain=475:540,       
    thick,     
    red,                    
] (
    {(1+0.25*cos(x/2)))*cos(x)},  % Punto en el borde (y = 0.5)
    {(1+0.25*cos(x/2)))*sin(x)},  
    {0.25*sin(x/2)}        
);

%------------------------------------- Circunferencia central e_3 ----------------------------------------
\addplot3 [
    samples=50,
    domain=232.5:-100,
    samples y=0,
    thick,
    violet,
    decoration={
        markings,
        mark=at position 0.9 with {\arrow{>}} 
    },
    postaction={decorate}
] (
    {cos(x)},
    {sin(x)},
    {0}
) node[pos=0.8, above] {$e_3$}; % <- El nombre se añade aquí

%-------------------------------------- Primer segmento frontera e_5 --------------------------------------------
\addplot3 [
    samples=50,
    samples y=0, % Evita saltos abruptos en y
    domain=180:-90,       
    thick,     
    blue,       
    init arrow             
] (
    {(1+0.25*cos(x/2)))*cos(x)},  % Punto en el borde (y = 0.5)
    {(1+0.25*cos(x/2)))*sin(x)},  
    {0.25*sin(x/2)}        
)node[pos=0.2, above] {$e_5$}; % <- El nombre se añade aquí
\addplot3 [
    samples=50,
    samples y=0, % Evita saltos abruptos en y
    domain=-151.75:-180,       
    thick,     
    blue,       
] (
    {(1+0.25*cos(x/2)))*cos(x)},  % Punto en el borde (y = 0.5)
    {(1+0.25*cos(x/2)))*sin(x)},  
    {0.25*sin(x/2)}        
);

%------------------------------- Corte transversal vertical superior e_1 ------------------------------------
\addplot3 [
    samples=50,
    domain=0:0.5,
    samples y=0,
    thick,
    lime,
    mid arrow,
] (
    {(1+0.5*x*cos(180/2)))*cos(180)}, % x: sustituye phi = 180
    {(1+0.5*x*cos(180/2)))*sin(180)}, % y: sustituye phi = 180
    {0.5*x*sin(180/2)}                % z: sustituye phi = 180
)node[pos=0.5, left] {$e_1$}; % <- El nombre se añade aquí

%------------------------------- Corte transversal vertical inferior e_2 ------------------------------------
\addplot3 [
    samples=50,
    domain=-0.5:0,
    samples y=0,
    thick,
    lime,
    mid arrow,
] (
    {(1+0.5*x*cos(180/2)))*cos(180)}, % x: sustituye phi = 180
    {(1+0.5*x*cos(180/2)))*sin(180)}, % y: sustituye phi = 180
    {0.5*x*sin(180/2)}                % z: sustituye phi = 180
)node[pos=0.5, left] {$e_2$}; % <- El nombre se añade aquí

%---------------------------------------------- Vértices ----------------------------------------------
% Vértices en phi = 180°
\node[right, yshift=5pt] at (axis cs:-1,0,-0.25) {$v_3$};
\node[right, yshift=5pt] at (axis cs:-1,0,0) {$v_2$};
\node[right, yshift=5pt] at (axis cs:-1,0,0.25) {$v_1$};
\addplot3[only marks, mark=*, mark size=1.5pt, black] coordinates {(-1,0,-0.25) (-1,0,0) (-1,0,0.25)};
\end{axis}
\end{tikzpicture}
\end{center}